\subsection{Módulos matemáticos}

Python incluye varios módulos que simplifican operaciones matemáticas frecuentes en programación competitiva.

El más utilizado es \texttt{math}.

\begin{lstlisting}[language=Python]
import math
\end{lstlisting}

\vspace{0.3cm}

\subsubsection{Redondeo}

\begin{center}
\begin{tabular}{|l|l|}
\hline
Función & Descripción \\ \hline
round(x,n) & Redondea a n decimales \\ \hline
math.floor(x) & Entero menor o igual \\ \hline
math.ceil(x) & Entero mayor o igual \\ \hline
math.trunc(x) & Elimina la parte decimal \\ \hline
\end{tabular}
\end{center}

\begin{lstlisting}[language=Python]
import math
math.floor(3.8) # 3
math.ceil(3.2)# 4
round(3.141592,2) # 3.14
math.trunc(7.98) # 7
factor = 10 ** N
floor_n = math.floor(x * factor) / factor
ceil_n  = math.ceil(x * factor) / factor
\end{lstlisting}

\vspace{0.3cm}

\subsubsection{Valor absoluto}

\begin{lstlisting}[language=Python]
abs(x)
math.fabs(x)
\end{lstlisting}

Normalmente basta utilizar

\begin{lstlisting}[language=Python]
abs(x)
\end{lstlisting}

\vspace{0.3cm}

\subsubsection{Raíces}

\begin{center}
\begin{tabular}{|l|l|}
\hline
Función & Uso \\ \hline
math.sqrt(x) & Raíz cuadrada \\ \hline
math.isqrt(x) & Raíz entera exacta \\ \hline
\end{tabular}
\end{center}

\begin{lstlisting}[language=Python]
math.sqrt(25) # 5.0
math.isqrt(30) # 5
\end{lstlisting}

En teoría de números suele preferirse \texttt{isqrt()}.

\vspace{0.3cm}

\subsubsection{Máximo común divisor y mínimo común múltiplo}

\begin{lstlisting}[language=Python]
math.gcd(a,b)

math.lcm(a,b)
\end{lstlisting}

También aceptan más de dos argumentos.

\begin{lstlisting}[language=Python]
math.gcd(a,b,c)

math.lcm(a,b,c)
\end{lstlisting}

Complejidad: \bigO{\log(\min(a,b))}

\vspace{0.3cm}

\subsubsection{Potencias}

\begin{lstlisting}[language=Python]
pow(a,b)
pow(a,b,m)
a ** b
math.pow(a,b)
\end{lstlisting}

Observaciones.

\begin{itemize}
\item \texttt{pow(a,b,m)} calcula la potencia módulo m de forma eficiente.
\item \texttt{**} devuelve enteros cuando es posible.
\item \texttt{math.pow()} siempre devuelve float.
\end{itemize}

\vspace{0.3cm}

\subsubsection{Factoriales}

\begin{lstlisting}[language=Python]
math.factorial(n)
\end{lstlisting}

Complejidad: \bigO{n}

\vspace{0.3cm}

\subsubsection{Combinatoria}

Desde Python 3.8 existen funciones incorporadas.

\begin{lstlisting}[language=Python]
math.comb(n,k)
math.perm(n,k)
\end{lstlisting}

Ejemplos.

\begin{lstlisting}[language=Python]
math.comb(5,2) # 10
math.perm(5,2) # 20
\end{lstlisting}

\vspace{0.3cm}

\subsubsection{Logaritmos}

\begin{lstlisting}[language=Python]
math.log(x)

math.log2(x)

math.log10(x)
\end{lstlisting}

\vspace{0.3cm}

\subsubsection{Trigonometría}

\begin{lstlisting}[language=Python]
math.sin(x)
math.cos(x)
math.tan(x)
math.asin(x)
math.acos(x)
math.atan(x)
\end{lstlisting}

Los ángulos siempre se expresan en radianes.

\vspace{0.3cm}

\subsubsection{Distancias}

\begin{lstlisting}[language=Python]
math.dist(p1,p2)

math.hypot(x,y)
\end{lstlisting}

\begin{lstlisting}[language=Python]
math.hypot(3,4) # 5
\end{lstlisting}

\vspace{0.3cm}

\subsubsection{Producto}

\begin{lstlisting}[language=Python]
math.prod(lista)
\end{lstlisting}

Ejemplo.

\begin{lstlisting}[language=Python]
math.prod([2,3,4]) # 24
\end{lstlisting}

\vspace{0.3cm}

\subsubsection{Constantes}

\begin{lstlisting}[language=Python]
math.pi
math.e
math.inf
math.nan
\end{lstlisting}

\vspace{0.3cm}

\subsubsection{Fracciones exactas}

Cuando los números en punto flotante producen errores de precisión puede utilizarse el módulo

\begin{lstlisting}[language=Python]
from fractions import Fraction
\end{lstlisting}

\begin{lstlisting}[language=Python]
Fraction(1,3)
Fraction(1,6)
Fraction(1,3)+Fraction(1,6)
\end{lstlisting}

Resultado.

\begin{lstlisting}[language=Python]
Fraction(1,2)
\end{lstlisting}

No existe error de redondeo.

\vspace{0.3cm}

\subsubsection{División entera}

En Python la división entera siempre redondea hacia menos infinito.

\begin{lstlisting}[language=Python]
7 // 2
# 3
-7 // 2
# -4
\end{lstlisting}

Este comportamiento es diferente al de C++ y Java.

\vspace{0.3cm}

\subsubsection{Aplicaciones frecuentes en ICPC}

Máximo común divisor.

\begin{lstlisting}[language=Python]
math.gcd(a,b)
\end{lstlisting}

Potencia modular.

\begin{lstlisting}[language=Python]
pow(a,b,MOD)
\end{lstlisting}

Raíz entera.

\begin{lstlisting}[language=Python]
math.isqrt(n)
\end{lstlisting}

Binomiales.

\begin{lstlisting}[language=Python]
math.comb(n,k)
\end{lstlisting}

Permutaciones.

\begin{lstlisting}[language=Python]
math.perm(n,k)
\end{lstlisting}

Fracciones exactas.

\begin{lstlisting}[language=Python]
Fraction(a,b)
\end{lstlisting}

\vspace{0.3cm}

\subsubsection*{Resumen}

\begin{center}
\begin{tabular}{|l|l|}
\hline
Función & Uso principal \\ \hline
floor() & Redondear hacia abajo \\ \hline
ceil() & Redondear hacia arriba \\ \hline
trunc() & Eliminar parte decimal \\ \hline
gcd() & Máximo común divisor \\ \hline
lcm() & Mínimo común múltiplo \\ \hline
sqrt() & Raíz cuadrada \\ \hline
isqrt() & Raíz entera \\ \hline
factorial() & Factorial \\ \hline
comb() & Combinaciones \\ \hline
perm() & Permutaciones \\ \hline
pow() & Potencias \\ \hline
log2() & Logaritmo base 2 \\ \hline
hypot() & Distancia euclidiana \\ \hline
prod() & Producto de un iterable \\ \hline
Fraction & Fracciones exactas \\ \hline
\end{tabular}
\end{center}